\chapter{Introduction} \label{ch:intro}

%\epigraph{``The most beautiful thing we can experience is the mysterious. It is the source of all true art and science.''}{Albert Einstein}
\epigraph{``What is the  gravitational field doing there, in such splendid isolation?''}{Bryce DeWitt}
\vspace{-0.5\baselineskip}

\bartriangledivide

\vspace{0.2\baselineskip}


This question posed by Bryce DeWitt in a text published posthumously in 2008
\cite{dewittQuantumGravityYesterday2009} strikes at the central issue that has
catalyzed the field we now call quantum gravity (QG). During the 20th century,
the physics community witnesses not one but two revolutions in our
understanding of the fundamental forces of nature. The development of general
relativity (GR) \cite{einsteinUberSpezielleUnd1956} radically changes our
understanding of the interaction we call gravity (and, with it, our
understanding of space and time). And at the same time the quantum revolution
starts
\cite{planckUberVerbesserungWienschen1900,planckUeberGesetzEnergieverteilung1901,einsteinUberErzeugungUnd1905,heisenbergUberQuantentheoretischeUmdeutung1925,bornZurQuantenmechanik1925,bornZurQuantenmechanikII1926,schrodingerUndulatoryTheoryMechanics1926}\footnote{Omissions
can not be avoided here.} and in the following years fundamentally alters our
picture of the nature of the atomic and subatomic and with it our understanding
of the remaining three interactions: electromagnetism, the weak interaction and
the strong interaction. Efforts to incorporate gravity have been undertaken
almost from the inception of quantum theory (QT), but the inclusion of the gravitational
interaction into the framework has met many obstacles. To this day, work on
formulating a comprehensive theory of quantum gravity remains ongoing and
continues to be a challenging pursuit. This is certainly what Bryce DeWitt is
referring to when he speaks about the ``splendid isolation'' of the
gravitational field. 
%For many people, the night sky is the first point of contact with the truly mysterious. For all of human existence, the stars in the sky have sparked curiosity and were the starting point of many scientific endeavors. Throughout human history, philosophers, scientists, and other scholars undertook many attempts to understand the position and movement of the stars in the sky. This academic pursuit culminates in the set of laws we surmise today under the name of ``gravity''. However, it is in the immediate human experience that binds us to the Earth; our understanding of it has enabled us instead to go beyond Earth. Launching a plethora of satellites into Earth's orbit and probes to other planets. It even enabled us to send humans to our most important satellite: the Moon.

Our first mathematical understanding of gravity came with the theory of
universal gravitation devised by Isaac Newton and first published in his
``Philosophiæ Naturalis Principia Mathematica'' (the mathematical principles of
natural philosophy) in 1678 \cite{newtonNewtonsPrincipiaMathematical1848}. This
theory should be superseded only after centuries by Albert Einstein's General
Theory of Relativity (GR) around 1916 \cite{einsteinUberSpezielleUnd1956}. It
marks a paradigm shift for modern physics. In the framework of GR, we can think
of gravity not as a typical force but as a result of the curvature of
spacetime; it is formulated as a purely geometrical theory (in the language of
differential geometry). This is summed up by the iconic phrase ``Spacetime
tells matter how to move; matter tells spacetime how to curve'' by John
Archibald Wheeler. 

%This fundamental difference in how we view gravity and spacetime is not the only break with the conventions of physics up to that point. While Newton's theory was devised from a large body of experimental data, the theory of GR is built upon principles, i.e., it is based on a set of expectations about the real world.

Moreover, there is an interesting epistemological difference between the two
theories. Newton's theory, although built on his principles of absolute space
and time, can not be directly derived from them. Rather, the inverse square law
was deduced from the experimental data available
\cite{cohenCambridgeCompanionNewton2002}.\footnote{The reference material
\cite{cohenCambridgeCompanionNewton2002} discusses the validity of this
statement with more nuance. However, it suffices here to clarify the difference
between approaches.} Einstein's theory of general relativity, on the other
hand, can be derived from first principles. Although there was experimental
evidence to build upon, namely the perihelion precession of Mercury and the
inherent tension between Newtonian theory and results from electromagnetism,
the theory was ultimately derived from these fundamentals rather than extracted
from the data. The principles guiding Einstein were Mach's principle, the
principle of equivalence, the principle of covariance, the principle of minimal
gravitational coupling, and the correspondence principle. Although researchers
agree on the fact that not all of these are necessary to derive general
relativity, there is no general agreement on the minimal set of principles
needed to derive the complete theory
\cite{dinvernoIntroducingEinsteinsRelativity2021}.

The second big revolution (quantum theory), on the other hand, is strongly
driven by accumulating experimental tensions. One prominent issue was termed
the ``ultraviolet catastrophe'' by Paul Ehrenfest
\cite{ehrenfestWelcheZugeLichtquantenhypothese1911}. This term describes the
inherent tension between the Rayleigh-Jeans law for radiation and experimental
observations, as well as the divergences it introduces
\cite{bartelmannTheoretischePhysik32018a}. Another prominent
example is the instability of the atom from a classical perspective. We know
from electromagnetism that a moving charge will emit radiation; classically,
this must then lead to a decay of the energy of the electrons in an atom.
Hence, the classical theory predicts that all electrons should, after some
time, ``fall'' into the nucleus, making atoms inherently unstable. Of course,
this is at odds with our daily experience of the existence of elements
\cite{bartelmannTheoretischePhysik32018a}. The development of
quantum theory finally resolved these and several other experimental tensions.
The results are startling even to the founders of QT: ``Everything is still
very vague and unclear to me, but it seems that electrons no longer move in
orbits.'' (Werner Heisenberg writing to Wolfgang Pauli)
\cite{polettiStrangenessQuantumMechanics2022}. Intense debates about the
physical meaning, the nature, and the interpretation of the new theory begin,
and many are still not resolved
\cite{jahnertRevisiting1927Solvay2012,cushingQuantumMechanicsHistorical1995}.

Nevertheless, work on quantum theory flourishes and soon leads to the
quantization of the electromagnetic interaction, establishing the framework of
Quantum Field Theory (QFT) \cite{diracQuantumTheoryEmission1997}. QFT unifies
quantum mechanics with special relativity by treating quantum fields, rather
than particles, as the fundamental objects of nature, where particles emerge as
localized excitations of these underlying fields. The foundational mathematical
structure of QFT is built upon symmetries, core principles that remain
unchanged under certain transformations. Lorentz invariance ensures that the
theory respects special relativity, meaning physical laws remain the same
regardless of the observer's relative motion. Meanwhile, gauge symmetries act
as local internal symmetries that dictate how particles can interact and what
charges they can carry. 


This first Quantum Field Theory of Quantum Electrodynamics (QED, Gauge Group
U(1)) remains to this day the most well-understood QFT of physical phenomena.
This initial success motivates the efforts to describe all fundamental forces
through QFT. After refining the methods and understanding of QFT through work
on QED, the weak force is tackled and finally quantized by introducing the
unified framework of Electroweak theory (Gauge Group U(1)$\times$SU(2))
\cite{glashowRenormalizabilityVectorMeson1959,salamWeakElectromagneticInteractions1959,weinbergModelLeptons1967}.
The strong force is quantized shortly thereafter through Quantum Chromodynamics
(Gauge Group
SU(3))\cite{grossUltravioletBehaviorNonAbelian1973,politzerReliablePerturbativeResults1973}.
Calculations in QCD are notoriously hard. It cannot be fully treated within the
framework of perturbation theory. This limitation arises from its distinct
behavior across energy scales, particularly in the low-energy regime where the
effective coupling constant becomes increasingly large. At these energies,
higher-order terms in the perturbative expansion become just as significant as
lower-order terms, making it impossible to treat the theory as small
corrections to a free theory. As we move to higher energies, however, these
theories exhibit a remarkable property called asymptotic freedom, where the
effective coupling approaches zero, making perturbative calculations
increasingly reliable. This was demonstrated in the seminal works of Gross,
Wilczek, and Politzer on asymptotic freedom
\cite{grossAsymptoticallyFreeGauge1973,davidpolitzerAsymptoticFreedomApproach1974}.

The methods of perturbation theory had been the dominant approach in studying
earlier QFT models and now had to be replaced by non-perturbative approaches
for studying the low energy regime of QCD. This fundamental issue was
recognized early in its development. Thus, the study of QCD has sparked
interest in many of the sophisticated tools for the non-perturbative study of
QFTs we use at present. Notable examples include functional Renormalization
Group methods and the lattice approach pioneered by Wilson
\cite{wilsonConfinementQuarks1974}. The lattice approach to quantum field
theory is essential to the theory we will direct our attention to in this
thesis. It is a reformulation where spacetime is discretized onto a
four-dimensional lattice, thereby providing a natural ultraviolet regulator
through the inverse lattice spacing. In this framework, gauge fields are
represented as group elements on the links of the lattice connecting
neighboring points, while the gauge field action is constructed from products
of these link variables around elementary squares (plaquettes) of the lattice
that ensure gauge invariance at finite lattice spacing. The quantum theory is
then formulated through a path integral where continuous spacetime integrals
become discrete sums \cite{munsterLatticeGaugeTheory2000}. This approach has
been very successful in studying QCD and other gauge theories
non-perturbatively.

After successfully quantizing three of the four known interactions of matter in
the universe with the methods of QFT, the fourth one was expected to follow
shortly thereafter. Instead, the fourth interaction turned out to be a lot more
resilient to any attempt to quantize it. Actually, work on this was started
quite early in the previously sketched timeline. In 1930 already, there was a
paper by Rosenfeld \cite{rosenfeldZurQuantelungWellenfelder1930} considering a
quantization of gravity.\footnote{This is usually referrenced as the starting
point of the field we know call quantum gravity. Still there are many earlier
works and discussions contributing to this development
\cite{peruzziTalesPrehistoryQuantum2018}.} The Chapel Hill conference on
gravity in 1957 presented a pivotal point in the history of the idea
\cite{ricklesRoleGravitationPhysics2011}. Here, Peter Bergmann, Bryce DeWitt,
Richard Feynman, John Archibald Wheeler, and many other influential physicists
of the time presented their arguments in favor of a quantum theory of gravity.
The discussion centered around issues like: What gravitational field is
produced by a quantum state in a superposition? What is the gravitational field
experienced by a quantum state in a superposition of location? Can spacetime be
in a superposition? These questions are still at the heart of the debate on the
quantization of gravity today.

There have been various arguments and thought experiments related to such
questions over the years
\cite{peresHybridClassicalquantumDynamics2001,dewittGlobalApproachQuantum2003,marlettoWhyWeNeed2017}
which have convinced a significant part of researchers to work in the field,
while some researchers remain skeptical of those arguments
\cite{oppenheimItTimeRethink2023}. As with any good debate in physics, this
will have to be settled by experimental evidence. Nevertheless, studying
possible candidates for a theory of quantum gravity is a worthwhile endeavor,
as it might point us to interesting phenomenology that might ultimately help in
the attempt to find real-world evidence for the need to quantize gravity.
Furthermore, this line of research has been very fruitful in further
understanding the established theories. 

For QFT, research into quantum gravity has yielded profound insights into the
theory's geometric and topological underpinnings. The topology of QFT was
studied in the context of topological quantum field theories (TQFT) \cite{atiyahTopologicalQuantumField1989}, where
physical measurements depend only on global, topological features of spacetime
rather than local geometric details. The geometric structure of gauge theories
refers to the elegant mathematical framework where gauge fields are understood
as connections on fiber bundles over spacetime, with gauge transformations
acting as geometric symmetries on these bundles. This perspective reveals how
Yang-Mills theory, for instance, can be interpreted through the lens of
differential geometry; its equations represent the curvature of connections on
principal bundles, while phenomena like instantons emerge as topologically
non-trivial solutions
\cite{ocampoGeometricTopologicalMethods2010,labastidaTopologicalQuantumField2005,nakaharaGeometryTopologyPhysics2018,naberTopologyGeometryGauge2011,hamiltonMathematicalGaugeTheory2017}.
Also, the study of quantum gravity renewed interest in and further refined the
non-perturbative methods for studying QFTs. 

Deeper insights into the structure of the theory were also gained in the case
of general relativity. A central insight concerns the nature of physical
observables, quantities that can be meaningfully measured in a
diffeomorphism-invariant theory. Since GR treats all coordinate systems as
equally valid (diffeomorphism invariance), true observables must be independent
of coordinate choices, which is not a trivial requirement
\cite{bergmannObservablesGeneralRelativity1961,perezObservablesQuantumGravity2010}.
Furthermore, various reformulations of GR have been studied. Through
reformulations like the (tetradic) Palatini formulation or the
MacDowell-Mansouri formulations, GR could be recast in a gauge-theoretic
framework. Reformulations like constrained $BF$ theory emphasize its relationship
to topological field theories
\cite{krasnovFormulationsGeneralRelativity2020}.
These alternative formulations are particularly valuable because they help
identify which degrees of freedom correspond to physical reality, as opposed to
mathematical redundancies in our description. Each reformulation provides a
different lens through which to view the physical content of the theory,
helping to separate the essential physical degrees of freedom from the
auxiliary mathematical structures we use to describe them. Fundamentally, this
work has deepened our appreciation of diffeomorphism invariance, not just as a
mathematical symmetry but as a profound physical principle. All of these
developments were in part driven and inspired by research into quantum gravity.

This monograph will be concerned with one of the many interesting attempts at
defining a consistent theory of QG. This attempt is called the spin foam
approach to quantum gravity.\footnote{A discussion of some other selected
approaches to quantum gravity and how they fit into the bigger space of
theories will be given in \cref*{sec:qg}.} The spin foam framework is part of
the larger class of functional integral (path integral) strategies for
quantizing gravity, sometimes called the sum over histories approach. This line
of research has been inspired by the success of the path integral formulation
of quantum mechanics and quantum field theory after its introduction to the
field by Richard Feynman \cite{feynmanSpaceTimeApproachNonRelativistic1948}. In
addition, the Spin Foam approach relies heavily on methods from Topological
Quantum Field Theory (TQFT) as defined by \cite{atiyahTopologicalQuantumField1989}.
In fact, a spin foam model of quantum gravity is defined by quantizing a $BF$
theory, which is a purely topological field theory, and subsequently imposing
so-called simplicity constraints. The simplicity constraints are what take the
purely topological $BF$ theory back to a description of gravity on the classical
level.\footnote{This process is discussed in more detail in
\cref{sec:bfgravity}.} In the following, we will shortly investigate the
development of our approach of choice: the Spin Foam approach.

We will take (without claiming uniqueness) the work of Turaev and Viro as the
starting point for the spin foam approach to quantum gravity
\cite{turaevStateSumInvariants1992}. As \cite{rovelliNotesBriefHistory2000}
states: ``Turaev and Viro define a state sum that on the one hand is a
rigorously defined TQFT, on the other hand can be seen as a regulated and well
defined version of the Ponzano-Regge quantization of 2+1 gravity.'' The work of
Ponzano and Regge \cite{ponanzaSemiclassicalLimitRacah1968} had already been
published 24 years in advance of the paper by Tuarev and Viro, but we want to
stress here the TQFT methods as a defining component for the spin foam
approach. Therefore, the work of Tuarev and Viro presents a better starting
point for the framework. The model in question has a 2+1 Riemannian signature
and is built upon the Group of SU(2). In addition to introducing the methods of
TQFT, Turaev and Viro used the quantum deformation of the group (SU(2)$_q$),
which amounts to a positive cosmological constant. The model of Tuarev and Viro
is a q-deformed version of the Ponanzo-Regge model. This introduces an upper
cut-off on representation labels, turning the diverging model of Ponzano and
Regge into a finite one. Shortly thereafter, Ooguri introduces a 4-dimensional
extension to the model of Ponanzo and Regge in
\cite{ooguriTopologicalLatticeModels1992}. He defines a topological model based
on a BF theory \cite{birminghamTopologicalFieldTheory1991} of the group SU(2)
in 4 dimensions. This model, now often called the Ooguri model, is still
essential in studying spin foam models for quantum gravity. It is one of the
first models that researchers will come in contact with when being introduced
to the spin foam approach to quantum gravity and is integral in the study of
some of the more advanced models. It is also the model that will be the main
focus of this monograph.

Shortly after this initial success, Reisenberger
\cite{reisenbergerWorldSheetFormulations1994}
was able to define a model for 4-dimensional self-dual Riemannian gravity in a
similar fashion by modifying the underlying TQFT. The same idea of building a
model based on a constrained TQFT\footnote{This is where the earlier mentioned
simplicity constraints come into play.} was then used by Barrett and Crane to
define a Riemannian model with a strong implementation of the constraints,
followed by a model with Lorentzian signature
\cite{barrettLorentzianSignatureModel2000}. Later this is succeeded by a model
developed by Engle, Pereira, Rovelli and Livine
\cite{engleLQGVertexFinite2008,livineConsistentlySolvingSimplicity2008,livineNewSpinfoamVertex2007}
parallel to Freidel and Krasnov \cite{freidelNewSpinFoam2008} which is
therefore sometimes called the EPRL-FK model. In this model, the boundary
states of the spin foam model connect to the states of Loop Quantum gravity.
This has shifted the focus of a large part of the Spin Foam community to the
EPRL-FK model.\footnote{This list is by no means exhaustive. Several models
that qualify as spin foam models for quantum gravity have been developed over
the years. For a more complete overview, the reader is referred to
\cite{perezSpinFoamApproach2012}.} 

To make sense of the state sum defined in any of these models (i.e., regularize
it), we need to introduce a discretization. A decomposition of the manifold
into a piecewise linear structure (also called a cell decomposition) can
accomplish that. Using this discretization, we can derive a form of the
partition function with variables and amplitudes associated to the dual complex
of this cell decomposition.\footnote{This is closely related to Lattice Gauge
Theories.} We will then have several integrals over representations that need
to be performed. This is, in general, an arduous task to perform numerically.
To make the expression more manageable for numerical computation, we introduce
a basis of intertwiners, turning the integrations into sums over intertwiner
labels. The computation of the state sum then translates to a tensor
contraction of large tensors associated to the vertices. The number of indices
of this tensor is dependent on the valency of the dual complex. The dimension
of these indices is represented by the intertwiner labels associated with the
edges and is determined by the spin representation labels associated with its
faces. Furthermore, there will be internal data (i.e., internal spin labels)
that we need to sum over. Explicitly evaluating such a state sum for the
aforementioned models is a complex task. At the time of writing, the only
analytical methods available can only be used for the asymptotic study of the
model.

The two limits that are most interesting for studying spin foam models are the
large spin behavior (asymptotic limit) and the large triangulation behavior
(refinement limit). In the limit of large spins, we expect the theory to
approach a semiclassical regime where the dominating contributions to the path
integral will be geometries that would correspond to geometries in Regge
Calculus, which is a discrete formulation of GR
\cite{mikovicEffectiveActionSemiclassical2011}. The large spin limit will lead
to a larger range of valid intertwiner labels, increasing the complexity of the
contraction between tensors. This means that the consumption of computational
resources will grow exponentially for the exact numerical evaluation methods.
The large triangulation limit can be used to investigate how the theory changes
under refinement. Studying such a refinement limit is one way to make sure that
we extract results from the theory that are independent of the regulator
\cite{asanteSpinFoamsRefinement2023}. Taking this limit will increase the
number of tensors appearing in the computation and additionally increase the
number of internal labels to be summed over. A large part of the numerical
efforts in the spin foam community are dedicated to investigating these two
limits.

Christensen and collaborators undertook a first attempt at the numerical study
of such models \cite{christensenNumericalEvidenceRegularized2009}, where
calculations were simplified to make the numerical evaluation possible. In this
paper the analysis was explicitly restricted to the one vertex case and
equilateral configurations only. Later, there was work on two vertices in the
Lorentzian case \cite{sarno2vertexLorentzianSpin2018}. The real breakthrough in
the numerical study came with the introduction of the \textsc{C++} library
\textsc{sl2cfoam} by Dona and collaborators
\cite{donaNumericalMethodsEPRL2018,donaNumericalStudyLorentzian2019,donaSearchingClassicalGeometries2020}.
Two previous developments were essential in making this breakthrough possible.
The first one concerned the reformulation of the computation of the 15j symbol
into a numerically more feasible form. First, in
\cite{donaSU2GraphInvariants2018}, it was decomposed into two 6j and one 9j
symbol, and then in \cite{donaNumericalMethodsEPRL2018}, further reduced into a
contraction of five 6j symbols. The second development were the Booster
matrices introduced by Speziale in \cite{spezialeBoostingWignersNjsymbols2017}.
The new formula for the 15j symbol greatly improved the calculation of the
SU(2) BF theory amplitude, and the Booster matrices of Speziale then made it
possible to define the EPRL-FK amplitude based on the SU(2) amplitude. This
toolkit was later remade and improved, creating the \textsc{sl2cfoam-next}
library \cite{donaSpinfoamsHighPerformanceComputing2023}, featuring a Julia interface,
suddenly making the numerical study of the EPRL-FK model and the Ooguri model
not only efficient and performant but also easier to set up.\footnote{We want
to direct the reader to \cite{donaHowToComputeEPRL2022}, where a complete
description of the necessary steps to set up a numerical evaluation of the
EPRL-FK model is given.} All of these developments have brought the study of
larger triangulations and large spins by explicit evaluation into reach for the
first time. Recently, there has been a proposal to improve the method of tensor
contractions by a cleverly chosen order of contractions to keep the tensors
involved in the computations as low-dimensional as possible
\cite{asanteEfficientTensorNetwork2024}. The authors show how this proposal can
drastically improve performance for the two previously mentioned models in the
presence of a reasonable amount of boundary data. 

Performant as these tools are, they still face the inherent problem of
exponential computational complexity. Recently, there have been efforts to
circumvent these problems entirely by mixing asymptotic methods with numerical
tools, giving rise to the study of ``Effective Spin Foams''
\cite{asanteEffectiveSpinFoam2020,asanteEffectiveSpinFoam2021,asanteDiscreteGravityDynamics2021}
and the method of complex critical points
\cite{hanAnalyticContinuationSpin2022,hanComplexCriticalPoints2022,hanComplexCriticalPoints2023}.
While both approaches already show promising results, they do not alleviate the
need to understand the full state sum, as they forgo the explicit calculation
of the vertex amplitude.

So, we must ask ourselves: Can we avoid the computational complexity of tensor
contractions and still evaluate the full model? We will have to add yet another
method to the computational toolkit for investigating spin foam models. A set
of methods that immediately comes to mind when faced with high computational
complexity are Monte Carlo (MC) methods, as they estimate the result with a
fixed number of samples. It remains an open question whether it is possible to
estimate the amplitude reasonably well using a number of samples smaller than
the total number of configurations. The quantities in question\footnote{By this
we mean the coherent vertex amplitude, a quantity that we will elaborate on in
\cref*{sec:su2cs}.} are highly oscillatory complex amplitudes and thus suffer
from the ``sign problem'' \cite{panSignProblemQuantum2024}. The problem refers
to the difficulty of sampling without missing the important contributions that
lead to the interference we expect in a quantum amplitude.\footnote{This issue
arises in complex oscillating amplitudes, but also in real amplitudes with
frequent sign changes.} Nevertheless, there have been examples of successful
applications of these methods in the spin foam literature.

In \cite{hanSpinfoamLefschetzThimble2021}, a Lefschetz thimble is defined by
the deformation of the integration contour, thereby fixing the imaginary part
of the amplitude to be constant on the integration contour. This allowed the
authors to avoid the sign problem when estimating the expectation values of
observables in the EPRL-FK model through MC. Remarkably, the authors of
\cite{donaSummingBulkQuantum2023} have shown that even without explicitly
avoiding the ``sign problem'' bulk intertwiners, as well as bulk spin degrees
of freedom, can be sampled using random sampling with a reasonable amount of
samples to estimate the amplitude of several relevant configurations. This
surprising result might suggest that the ``sign problem'' is not as severe in
the studied spin foam models. Building on this success, the author of
\cite{steinhausMonteCarloAlgorithm2024} shows that coherent boundary
intertwiners can be sampled as well based on the coherent state associated to
that amplitude. This is implemented for one vertex in the Ooguri model. 

This thesis builds on the pioneering work of Steinhaus. We have developed a
novel approach to bulk intertwiner sampling that draws inspiration from
effective spin foam models, employing classical geometric configurations
derived from area-length relationships to inform Monte Carlo proposal
distributions. While this represents a meaningful step forward in our
computational toolkit, it is important to acknowledge that further
investigations will be necessary to fully determine the potential of these
methods.  

Our geometrically-motivated sampling strategy represents one of the first
systematic attempts to leverage semiclassical intuition for importance sampling
in the quantum regime of spin foam models. The preliminary results suggest that
this approach may offer computational advantages over random sampling in
certain regimes, particularly when Euclidean configurations dominate
the quantum amplitudes as theoretical expectations suggest.  

The work is not solely focused on computational efficiency gains. Our methods
provide new tools for investigating fundamental questions about the
quantum-to-classical transition in discrete gravity. Due to its setup the
method lends itself to probing the geometries that contribute to the path
integral. In taking inspiration from the effective spin foam approach it
presents a tool to understand how well geometrical constraints are implemented
in the quantum model.  

The techniques we develop for efficient computation of SU(2) $BF$ amplitudes can
be straightforwardly extended to more realistic models, with the EPRL/FK model
being a natural next target given its heavy reliance on SU(2) $BF$ structures
by construction. By
establishing these foundational methods and demonstrating their viability in
the simpler topological context, this work provides essential stepping stones
toward tackling the more complex and physically relevant spin foam models. 

%The ``informed'' guessing
%will entail finding the classical geometric configurations that are consistent
%with a given spin configuration and then using the methods of
%\cite{steinhausMonteCarloAlgorithm2024} to create the sampling distributions.
%We will investigate the structure of bulk intertwiners for some examples and,
%from there, develop the method.

%We want to highlight one of these arguments by Richard Feynman \cite{feynmanFeynmanLecturesGravitation1995}. 

% \begin{figure}[h]
%     \centering
%     \resizebox{85mm}{!}{\input{Pics/drawings/QGdoubleslit.pdf_tex}}
%     \caption{The double slit experiment for gravity as presented in \cite{feynmanFeynmanLecturesGravitation1995}. Electrons are sent through the experiment and will be found to be in a superposition of the state associated with the upper and the state associated with the lower slit. An interference is observed at the detector.}
%     \label{fig:qgdoubleslit}
% \end{figure}


% In \cref{fig:qgdoubleslit}, the general idea is presented. We consider a modified version of the famous double slit experiment such that there is a mass present directly behind the divider with the slits. That mass will experience the (admittedly tiny) gravitational field produced by the electron from the usual double-slit setup. We ask ourselves: What is the gravitational field experienced by this mass? In the Quantum Mechanical picture, the electrons are in a superposition of going through the upper and lower slit $\frac{1}{\sqrt{2}}(\ket{U}+\ket{L})$; therefore, the gravitational field experienced by the mass needs to be in a superposition as well, or the typical interference pattern will not be observed at the detector. We know that this interference pattern is actually observed in the usual double-slit experiments \cite{joenssonElektroneninterferenzenMehrerenKuenstlich1961}, and as far as we know, the gravitational interaction can not be screened. The fact that the gravitational interaction cannot be screened would mean that the mass in \cref{fig:qgdoubleslit} is just included for illustration but is not actually needed as the whole experimental setup will already experience the questionable gravitational field. Thus, only possibility one is logically permitted.\footnote{Some interpretations of Quantum theory might allow to avoid this conclusion. A notable example would be the interpretation of many worlds, where the mass would experience different gravitational fields in different worlds. As with any thought experiment, the implications are heavily debated in the literature.} At least, this is how the argument is usually presented.
